\section{结论}

在所测试的 OOD 场景中，许多 VLA 失败表现为感知、语言或协调局限，而不只是运动能力不足。
为解决这一问题，我们提出 \sys：一种记忆增强策略外壳。它通过动态 MCP 工具编排以及由检索、归因驱动推理和双阶段记忆巩固组成的闭环来包裹冻结的 VLA 策略，无需进行任何参数微调。
在多个 VLA 骨干模型上的评估表明，\sys\ 在多种 OOD 条件下都能持续改善上下文内自适应和鲁棒性。
未来工作将扩展 MCP 工具库以覆盖更广泛的失败模式，研究前沿多模态模型以进一步细化失败归因粒度，并探索这类记忆引导的干预如何迁移到更多样的真实机器人平台。

\section*{致谢}

我们衷心感谢匿名审稿人；他们的评审意见、反馈和建议显著改进了本工作。
本工作部分受到上海市科学技术委员会新一代信息技术项目（编号 25511104100）、国家自然科学基金（编号 62432010）以及教育部基础学科和交叉学科突破计划（JYB2025XDXM122）支持。
本工作还受到 OpenHarmony 具身智能项目管理委员会（PMC）支持。

\endinput
